\chapter{Imagen espectral--polar y reducción exacta de Riemann--Weil}

\section{Punto de partida: la estructura espectral--polar completa}

La entrada de la resolución no es la forma de Weil aislada ni una lista de
ceros. Es la estructura espectral--polar completa producida por el constructor
común y acreditada por el propietario exacto que precede a este capítulo; no es
un dato externo. La cadena causal propia de este capítulo es
\begin{equation}\label{eq:pdf2-rh-cadena-procedencia}
 \mathfrak G_{\rm sp}
 \xrightarrow{\ \operatorname{pr}_{X,{\rm sp}}\ }X_{\chi_{\rm sp}}
 \xrightarrow{\ \mathcal A_{\rm sp}\ }\mathfrak M_{\rm sp}
 \xrightarrow{\ \mathcal P_{\rm sp}\ }\mathcal C_{\rm GW}.
\end{equation}
El codominio clásico aparece sólo en la última flecha; no selecciona
retrospectivamente el estado, las cartas, la proyección ni el terminal.

La estructura de partida contiene el objeto de trece coordenadas
\begin{equation}\label{eq:pdf2-rh-res-sp-completa}
\begin{aligned}
 \mathfrak G_{\rm sp}^{(13)}=\bigl(&
 \mathcal F_{\rm sp},
 \operatorname{Ext}_{\Gamma,{\rm sp}},
 \operatorname{Rel}_{\rm sp},
 \mathcal N_{\rm sp},
 \operatorname{or}_{\rm sp},
 \operatorname{Carr}_{\rm sp},
 \operatorname{Mem}_{\rm sp},\\
 &\partial_{\rm sp},
 \gamma_{\rm sp},
 \Sigma_{\rm sp}^{\rm mult},
 \operatorname{Surv}_{\rm sp},
 \tau_{\rm sp}^{\rm term},
 \mathcal L_{\rm sp}\bigr).
\end{aligned}
\end{equation}
Cada coordenada interviene materialmente:
\begin{enumerate}
 \item \(\mathcal F_{\rm sp}\) conserva las dos hojas, su orientación, el
 par residuo--cociente y la firma de la publicación doble.
 \item \(\operatorname{Ext}_{\Gamma,{\rm sp}}\) transporta esa fibra por la
 clase dirigida de cortes y niveles nonádicos.
    \item \(\operatorname{Rel}_{\rm sp}\) contiene el cociclo de acarreo, la
    ortogonalidad \(P=P^*=P^2\), las dos publicaciones y el complemento
    estructural \(B_{\rm sp}^{\rm str}:=\Xi^*(I-P)\Xi\). La identificación
    de este complemento con la forma de Guinand--Weil es una flecha distinta
    y se audita por separado en este capítulo.
 \item \(\mathcal N_{\rm sp}\) registra el nivel \(m\), la ventana \(R\),
 el cociente \(q=5/9\), los momentos de ambas hojas y los índices de los
 canales primo, gamma y polar.
 \item \(\operatorname{or}_{\rm sp}=\widehat R=(R,K)\) mantiene distintas
 las dos cartas orientadas.
 \item \(\operatorname{Carr}_{\rm sp}=\operatorname{Carr}_R\) es el
 acarreo operatorial y no se anula al publicar una traslación.
 \item \(\operatorname{Mem}_{\rm sp}=(S_0,M,\mathcal O_9)\) conserva la
 identidad finita de Stein y su terminal.
 \item \(\partial_{\rm sp}=(I-P)\Xi\) es la frontera que no aparece en la
 compresión \(P\Xi\) y cuya energía define el complemento estructural. Su
 identificación con la forma de Weil no se incorpora a esta definición.
 \item \(\gamma_{\rm sp}\) es la ruta cofinal de ventanas, con un único
 regulador para todos los canales.
 \item \(\Sigma_{\rm sp}^{\rm mult}\) es la multisección de las dos
 publicaciones y sus componentes primo--gamma--polar; no es la firma de un
 único estado.
 \item \(\operatorname{Surv}_{\rm sp}\) conserva la familia bajo
 \(m\mapsto m+1\) y \(\mathcal D_R\hookrightarrow\mathcal D_{R'}\).
 \item \(\tau_{\rm sp}^{\rm term}\) conserva el terminal
 \(M^{*N}S_0M^N\), que en la carta espejo se publica mediante \(E_R\).
    \item \(\mathcal L_{\rm sp}=(\Xi,P,B_{\rm sp}^{\rm str})\) es el lector final; se aplica
 después de antecedente, cartas, memoria, ruta y terminal.
\end{enumerate}

El ensamblador y el publicador son
\begin{equation}\label{eq:pdf2-rh-ensamblador-sp}
\begin{aligned}
 \mathcal A_{\rm sp}(\operatorname{pr}_{X,{\rm sp}}\widehat x)
  ={}&(\widehat X_\infty,\Gamma_9,\widehat R,
      \operatorname{Carr}_R,S_0,M,\mathcal O_9,
      \Xi,P,\mathcal D_R^{\rm cof}),\\
 \mathfrak M_{\rm sp}={}&\operatorname{Im}\mathcal A_{\rm sp},\\
 \mathcal P_{\rm sp}(\mathfrak M_{\rm sp})
  ={}&(B_{\rm sp}^{\rm str},\Xi^*(I-P)\Xi,
       \mathscr I_{\rm GW},\Delta_{\rm bind},
       \text{familia cofinal primo--gamma--polar}),\\
 \Delta_{\rm bind}:={}&B_{\rm sp}^{\rm str}-\mathscr I_{\rm GW}.
\end{aligned}
\end{equation}
La memoria finita del objeto de partida satisface
\begin{equation}\label{eq:pdf2-rh-stein-finito}
 S_0=\sum_{r=0}^{N-1}M^{*r}\mathcal O_9^*\mathcal O_9M^r
     +M^{*N}S_0M^N.
\end{equation}
El último sumando permanece visible hasta que se demuestra su extinción. Por
tanto, reemplazar \(\mathfrak G_{\rm sp}\) por la sombra publicada
\(B_{\rm sp}^{\rm str}\) cambia el objeto de partida; no es una abreviación
del mismo teorema. El residual \(\Delta_{\rm bind}\) se conserva hasta que la
coligación source--first demuestra materialmente su anulación en
\eqref{eq:pdf2-rh-binding-source-first}.

\section{Dominio cofinal y forma objetivo}

Para $R>0$ se fija
\[
 \mathcal D_R=C_c^\infty((-R/2,R/2)),\qquad
 \widehat\psi(t)=\int_{\mathbb R}\psi(x)e^{-itx}\,dx,
\]
y, para $\psi,\phi\in\mathcal D_R$,
\[
 h_{\psi,\phi}(z)=\widehat\psi(z)
 \overline{\widehat\phi(\overline z)},\qquad
 \check h_{\psi,\phi}(a)=\langle\psi,T_a\phi\rangle.
\]
La unión $\bigcup_R\mathcal D_R=C_c^\infty(\mathbb R)$ es exacta. Para todo
conjunto finito de puntos complejos distintos, las evaluaciones
$\psi\mapsto(\widehat\psi(z_1),\ldots,\widehat\psi(z_n))$ son conjuntamente
sobreyectivas cuando $R$ es suficientemente grande. Por tanto, esta clase cofinal
separa las evaluaciones requeridas por el criterio de Weil.

\begin{lemma}[Prueba de la separación conjunta]
\label{lem:pdf2-rh-separacion}
La aplicación
\[
 \mathcal D_R\longrightarrow\mathbb C^n,
 \qquad
 \psi\longmapsto
 (\widehat\psi(z_1),\ldots,\widehat\psi(z_n))
\]
es sobreyectiva para todo conjunto de puntos complejos distintos
\(z_1,\ldots,z_n\), una vez que \(R\) contiene un intervalo no vacío.
\end{lemma}

\begin{proof}
Si los funcionales de evaluación fueran linealmente dependientes, existirían
\(c_1,\ldots,c_n\), no todos nulos, tales que
\[
 \int_{\mathbb R}\psi(x)
 \left(\sum_{j=1}^nc_je^{-iz_jx}\right)\,dx=0
 \qquad(\psi\in\mathcal D_R).
\]
La función analítica \(\sum_jc_je^{-iz_jx}\) se anularía en
\((-R/2,R/2)\), luego en todo el plano. La independencia de exponenciales
con exponentes distintos fuerza \(c_j=0\) para todo \(j\), contradicción.
Los funcionales son independientes; como el codominio es finito-dimensional,
la aplicación tiene rango \(\mathbb C^n\).
\end{proof}

\section{Las dos columnas y los tres canales calculados antes del signo}

Sea $\vartheta_L$ un regulador común. Con
\[
 \ell_{p,k}=k\log p,\qquad w_{p,k}=(\log p)p^{-k/2},\qquad
 \mu_\Gamma(a)=\frac{e^{-a/2}}{1-e^{-2a}},
\]
se usan la inclusión uniforme $j_m$ y el canal de costura
$\widehat\delta_{a,m}$, que satisfacen
\[
 j_m^*j_m=I,\qquad
 \widehat\delta_{a,m}^*\widehat\delta_{b,m}
 =(I-T_a)^*(I-T_b).
\]
Las columnas reguladas son
\begin{align*}
 J_{+,m,R,L}\psi={}&
 \Bigl(
  (\sqrt{\vartheta_L(\ell_{p,k})w_{p,k}}\,
       \widehat\delta_{\ell_{p,k},m}\psi)_{\ell_{p,k}\leq R},\\
 &\hspace{9mm}a\mapsto
  \sqrt{\vartheta_L(a)\mu_\Gamma(a)}\,
       \widehat\delta_{a,m}\psi,\ \zeta_+b_+(\psi)
 \Bigr),\\
 J_{-,m,R,L}\psi={}&
 \Bigl(
  (\sqrt{2\vartheta_L(\ell_{p,k})w_{p,k}}\,j_m\psi)_{\ell_{p,k}\leq R},
  \sqrt{-c_\Gamma}\,j_m\psi,\ \zeta_-b_-(\psi)
 \Bigr),
\end{align*}
donde $c_\Gamma=\psi_\Gamma(1/4)-\log\pi<0$ y
$b_\pm=(\widehat\psi(i/2)\pm\widehat\psi(-i/2))/\sqrt2$.

Una expansión término a término, sin usar ceros ni el signo buscado, da
\begin{align*}
 &\langle J_{+,m,R,L}\psi,J_{+,m,R,L}\phi\rangle
 -\langle J_{-,m,R,L}\psi,J_{-,m,R,L}\phi\rangle
 =\mathscr I_{{\rm GW},L}(\psi,\phi),\\
 \mathscr I_{{\rm GW},L}(\psi,\phi)
 ={}&h_{\psi,\phi}(i/2)+h_{\psi,\phi}(-i/2)
 +c_\Gamma\check h_{\psi,\phi}(0)\\
 &+\int_0^\infty\vartheta_L(a)\mu_\Gamma(a)
 [2\check h(0)-\check h(a)-\check h(-a)]\,da\\
 &-\sum_{p,k}\vartheta_L(\ell_{p,k})w_{p,k}
 [\check h(\ell_{p,k})+\check h(-\ell_{p,k})].
\end{align*}
El truncamiento de ambas columnas en \(\ell_{p,k}\leq R\) es obligatorio, no
notacional. Si \(\psi,\phi\in\mathcal D_R\), entonces
\(\check h_{\psi,\phi}(\pm\ell)=0\) para \(\ell>R\); la contribución de cada
par de coordenadas omitido a la \emph{diferencia} de Grams es por tanto cero.
La columna negativa no truncada, en cambio, tendría norma infinita porque
\(\sum_p(\log p)p^{-1/2}=\infty\). En el conjunto finito
\(\ell_{p,k}\leq R\) las coordenadas primas convergen individualmente cuando
\(L\to\infty\). En el canal gamma, el integrando es \(O(a)\) en cero y
\(\mu_\Gamma(a)=O(e^{-a/2})\) en infinito, de modo que la convergencia también
es en norma. Definimos, por tanto,
\begin{equation}\label{eq:pdf2-rh-columnas-reducidas}
 J_{\pm,m,R}:=\lim_{L\to\infty}J_{\pm,m,R,L}
 \quad\text{en sus Hilbert de coordenadas truncadas},
\end{equation}
y suprimimos \(m\) cuando permanece fijo. Por convergencia dominada,
\[
 \mathscr I_{{\rm GW},L}\longrightarrow\mathscr I_{\rm GW}
 \qquad(L\to\infty),
\]
y las identidades son naturales bajo $m\mapsto m+1$ y $R\le R'$.

Las columnas reducidas \(J_{+,R}\) y \(J_{-,R}\) son ahora operadores
materiales de \(\mathcal D_R\) a Hilbert. La expansión explícita prueba, sin
consultar ceros ni signo,
\begin{equation}\label{eq:pdf2-rh-gram-difference-gw}
 \langle J_{+,R}f,J_{+,R}g\rangle
 -\langle J_{-,R}f,J_{-,R}g\rangle
 =\mathscr I_{\rm pr}(f,g)+\mathscr I_\Gamma(f,g)
  +\mathscr I_{\rm pol}(f,g)
 =\mathscr I_{\rm GW}(f,g).
\end{equation}
Los relojes \(\ell_{p,k}\) producen el sumando primo; la integral con
\(\mu_\Gamma\) y \(c_\Gamma\check h(0)\), el sumando gamma; y
\[
 \langle b_+(f),b_+(g)\rangle-
 \langle b_-(f),b_-(g)\rangle
 =h_{f,g}(i/2)+h_{f,g}(-i/2)
\]
produce el sumando polar. Esta parte del cálculo es material e independiente
del puente estructural.

\section{Coligación source--first de la imagen espectral--polar}

Fijamos primero un nivel nonádico \(m\); este índice se suprime en
\(J_{\pm,R}\), \(E_R\) y \(\Sigma_R\) hasta construir el codificador. La
descomposición de los codominios explícitos de
\eqref{eq:pdf2-rh-columnas-reducidas} es
\[
 I_+=\{\mathrm{pr},\Gamma,\mathrm{pol}\},\qquad
 I_-(R)=\{\Gamma0,\mathrm{pol}\}
 \sqcup\{(p,k):\ell_{p,k}\leq R\},
\]
\[
 \mathscr H_{+,R}^{\rm amb}=\bigoplus_{i\in I_+}H_{+,i,R},
 \qquad
 \mathscr H_{-,R}^{\rm amb}=\bigoplus_{j\in I_-(R)}H_{-,j,R}.
\]
No identificamos un subespacio alcanzable con una suma de sus proyecciones
coordenadas. Definimos, con las inclusiones ambientales entendidas,
\[
 H_{+,R}:=\overline{\operatorname{Ran}J_{+,R}},\qquad
 H_{-,R}:=\overline{\operatorname{Ran}J_{-,R}},
\]
como subespacios cerrados de los dos Hilbert ambientales. El propietario
espectral impreso antes de este capítulo construye
\(\mathcal H_R^{\rm sh}\), la diferencia antiespecular y la fila isométrica
source--first mediante
\eqref{eq:amp-rh-owner-q}--\eqref{eq:amp-rh-owner-binding}, prueba su
telescopía en \eqref{eq:amp-rh-owner-telescopia} y publica la extinción
terminal en \cref{thm:amp-rh-cierre-cofinal-propietario}. La realización
que sigue conserva esa construcción, antes de la identificación clásica,
en la fila
\begin{equation}\label{eq:pdf2-rh-coligacion-source-first}
 \Sigma_R=
 \begin{pmatrix}\mathcal K_R^{\rm src}\\ \mathcal F_R^{\rm src}\end{pmatrix}:
 H_{+,R}\longrightarrow
 H_{-,R}\oplus\mathcal H_R^{\rm sh},
 \qquad
 \Sigma_R^*\Sigma_R
 = (\mathcal K_R^{\rm src})^*\mathcal K_R^{\rm src}
  +(\mathcal F_R^{\rm src})^*\mathcal F_R^{\rm src}=I_{H_{+,R}}.
\end{equation}
Su acción alcanzable se fija sobre la columna positiva completa por
las relaciones de las dos publicaciones que forman parte de
\(\operatorname{Rel}_{\rm sp}\). Para evitar cualquier sustitución de lector,
denotamos por
\((\mathcal X_R^{\rm str},\Xi_R^{\rm str},P_R^{\rm str})\) la restricción a
\(\mathcal D_R\) del lector estructural de
\eqref{eq:pdf2-rh-res-sp-completa}. Sus dos momentos son, literalmente,
\begin{equation}\label{eq:pdf2-rh-momentos-lector-estructural}
\begin{aligned}
 \langle\Xi_R^{\rm str}f,\Xi_R^{\rm str}g\rangle
  &=\langle J_{+,R}f,J_{+,R}g\rangle,\\
 \langle P_R^{\rm str}\Xi_R^{\rm str}f,
          P_R^{\rm str}\Xi_R^{\rm str}g\rangle
  &=\langle J_{-,R}f,J_{-,R}g\rangle,
 \qquad (P_R^{\rm str})^*= (P_R^{\rm str})^2=P_R^{\rm str}.
\end{aligned}
\end{equation}
La coordenada de frontera viene con la isometría estructural
\[
 \jmath_R^{\rm sh}:
 \overline{\operatorname{Ran}((I-P_R^{\rm str})\Xi_R^{\rm str})}
 \longrightarrow\mathcal H_R^{\rm sh}.
\]
Por definición de la fila source--first, no por una identificación posterior,
\begin{equation}\label{eq:pdf2-rh-source-state-output}
 E_R:=\mathcal F_R^{\rm src}J_{+,R}
     =\jmath_R^{\rm sh}(I-P_R^{\rm str})\Xi_R^{\rm str}:
 \mathcal D_R\longrightarrow\mathcal H_R^{\rm sh},
 \qquad
 \Omega_{0,R}:=\mathcal K_R^{\rm src}J_{+,R}-J_{-,R}.
\end{equation}
No hay una coligación por canal: la misma fila recibe simultáneamente los
relojes primo--potencia, la integral gamma, el modo escalar y las dos líneas
polares. En particular, conserva todos los productos cruzados que aparecen al
polarizar \eqref{eq:pdf2-rh-gram-difference-gw}.
Con las proyecciones ambientales \(\pi_j^-:\mathscr H_{-,R}^{\rm amb}
\to H_{-,j,R}\), las filas de salida correctamente tipadas son
\begin{equation}\label{eq:pdf2-rh-bloques-coligacion}
 \mathcal K_{j,R}^{\rm src}
 :=\pi_j^-\mathcal K_R^{\rm src}:
 H_{+,R}\longrightarrow H_{-,j,R},
 \qquad
 \mathcal F_R^{\rm src}:H_{+,R}\longrightarrow\mathcal H_R^{\rm sh}.
\end{equation}
Así no se aplica \(\mathcal K_R^{\rm src}\) a una coordenada positiva que
pudiera no pertenecer por sí sola a \(H_{+,R}\). Todos los cruces entre
canales permanecen dentro de
\((\mathcal K_R^{\rm src})^*\mathcal K_R^{\rm src}\) y
\((\mathcal F_R^{\rm src})^*\mathcal F_R^{\rm src}\); no se ortogonalizan
antes de la identidad de energía.

La coordenada de supervivencia construida por
\eqref{eq:amp-rh-owner-q}, \eqref{eq:amp-rh-owner-telescopia} y
\eqref{eq:amp-rh-owner-terminal} produce los espacios residuales
\(\mathcal R_{n,R}^{\rm sh}\), continuaciones
\(\Omega_{n,R}:\mathcal D_R\to\mathcal R_{n,R}^{\rm sh}\) e isometrías
antiespeculares
\(V_{n,R}:\mathcal R_{n+1,R}^{\rm sh}\to\mathcal R_{n,R}^{\rm sh}\). Su
primer espacio y su primer operador son, por definición tipada,
\[
 \mathcal R_{0,R}^{\rm sh}:=H_{-,R},\qquad
 \Omega_{0,R}:=\mathcal K_R^{\rm src}J_{+,R}-J_{-,R}.
\]
Su recurrencia de una vuelta completa es la publicación cofinal del mismo
operador construido, no un dato añadido a la realización:
\begin{equation}\label{eq:pdf2-rh-recurrencia-residual}
 q=\frac59,\qquad
 \Omega_{n,R}=qV_{n,R}\Omega_{n+1,R},
 \qquad V_{n,R}^*V_{n,R}=I,
 \qquad
 \sup_{n\geq0}\|\Omega_{n,R}f\|<\infty
 \quad(f\in\mathcal D_R).
\end{equation}
Esta recurrencia pertenece a la coligación enriquecida: transporta la
diferencia de salida, no la forma de Weil ni su signo.

\begin{proposition}[Anulación coinductiva de la diferencia de salida]
\label{prop:pdf2-rh-omega-cero}
Para todo \(R>0\) y todo \(f\in\mathcal D_R\),
\begin{equation}\label{eq:pdf2-rh-output-binding-source}
 \Omega_{n,R}f=0\quad(n\geq0),
 \qquad
 \boxed{\mathcal K_R^{\rm src}J_{+,R}=J_{-,R}}.
\end{equation}
\end{proposition}

\begin{proof}
Iterar \eqref{eq:pdf2-rh-recurrencia-residual} \(N\) veces y usar que cada
\(V_{n,R}\) es isométrica da
\[
 \|\Omega_{n,R}f\|
 =q^N\|\Omega_{n+N,R}f\|
 \leq q^N\sup_{j\geq n}\|\Omega_{j,R}f\|.
\]
El miembro derecho converge a cero porque \(0<q=5/9<1\). Por tanto
\(\Omega_{n,R}f=0\); el caso \(n=0\) y
\eqref{eq:pdf2-rh-source-state-output} producen la identidad enmarcada.
\end{proof}

La acción fila a fila de la salida ya no queda implícita. Las proyecciones son
las restricciones a \(H_{-,R}\) de las proyecciones del Hilbert ambiental, y
\begin{equation}\label{eq:pdf2-rh-salidas-generador-a-generador}
\begin{aligned}
 \mathcal K_{\Gamma0,R}^{\rm src}J_{+,R}f
  &=\sqrt{-c_\Gamma}\,j_mf,\\
 \mathcal K_{(p,k),R}^{\rm src}J_{+,R}f
  &=\sqrt{2w_{p,k}}\,j_mf
  &&(\ell_{p,k}\leq R),\\
 \mathcal K_{{\rm pol},R}^{\rm src}J_{+,R}f
  &=\zeta_-b_-(f).
\end{aligned}
\end{equation}
Estas filas agotan \(J_{-,R}\), mientras
\begin{equation}\label{eq:pdf2-rh-frontera-generador-a-generador}
 \mathcal F_R^{\rm src}J_{+,R}f=E_Rf
\end{equation}
conserva, con todos sus cruces, la coordenada de frontera.
Evaluar \eqref{eq:pdf2-rh-coligacion-source-first} en \(J_{+,R}f\) y
\(J_{+,R}g\), y usar \eqref{eq:pdf2-rh-output-binding-source}, produce la
identidad de energía anterior al signo
\begin{equation}\label{eq:pdf2-rh-energia-source-first}
 \boxed{
 \langle J_{+,R}f,J_{+,R}g\rangle
 =\langle J_{-,R}f,J_{-,R}g\rangle
  +\langle E_Rf,E_Rg\rangle.}
\end{equation}
Al compararla con el cálculo independiente
\eqref{eq:pdf2-rh-gram-difference-gw} se obtiene, sin introducir una raíz de
la forma objetivo,
\begin{equation}\label{eq:pdf2-rh-binding-source-first}
 \boxed{E_R^*E_R=\mathscr I_{\rm GW}\big|_{\mathcal D_R}.}
\end{equation}
En particular,
\(\Delta_{{\rm bind},R}
:=E_R^*E_R-\mathscr I_{\rm GW}|_{\mathcal D_R}=0\); esta anulación es la
salida de \eqref{eq:pdf2-rh-coligacion-source-first}--
\eqref{eq:pdf2-rh-energia-source-first}, no una premisa insertada en el
publicador.

\section{Complemento estructural, hoja espejo y telescopía terminal}

El estado \(E_R=\mathcal F_R^{\rm src}J_{+,R}\) construido por la fila
source--first es la coordenada de frontera de
\(\mathfrak G_{\rm sp}\). Definimos su forma estructural por
\begin{equation}\label{eq:pdf2-rh-bw-estructural}
 \mathcal B_{{\rm sp},R}^{\rm str}
 :=E_R^*E_R
 =\mathscr I_{\rm GW}\big|_{\mathcal D_R},
\end{equation}
donde la segunda igualdad es
\eqref{eq:pdf2-rh-binding-source-first}, ya deducida de la coligación y de
la expansión generador a generador.

Sean \(P_+^{\rm sh}\) y \(P_-^{\rm sh}\) los proyectores ortogonales
complementarios de las dos hojas. Se fijan
\begin{equation}\label{eq:pdf2-rh-q-a-d}
 q=\frac59,\qquad
 A=P_+^{\rm sh}+qP_-^{\rm sh},\qquad
 D=\sqrt{1-q^2}\,P_-^{\rm sh}.
\end{equation}
La ortogonalidad da la identidad exacta
\begin{equation}\label{eq:pdf2-rh-conservacion-local}
 A^*A+D^*D
 =P_+^{\rm sh}+q^2P_-^{\rm sh}
 +(1-q^2)P_-^{\rm sh}=I.
\end{equation}
Las nueve fases coordinan una sola vuelta con memoria:
\begin{equation}\label{eq:pdf2-rh-monodromia}
 M_9=\Phi A_8\cdots A_0=qV_9,
 \qquad V_9^*V_9=I.
\end{equation}
La contracción por vuelta es \(q\), no \(q^9\); las fases internas conservan
el orden y el acarreo.

La coligación tipa
\(E_R:\mathcal D_R\to\operatorname{Ran}P_-^{\rm sh}\). Por tanto,
\begin{equation}\label{eq:pdf2-rh-binding-espejo}
 P_+^{\rm sh}E_R=0,
 \qquad AE_R=qE_R,
 \qquad
 E_R^*E_R=\mathcal B_{{\rm sp},R}^{\rm str}
 =\mathscr I_{\rm GW}\big|_{\mathcal D_R}.
\end{equation}

\begin{proposition}[Telescopía finita con terminal conservado]
\label{prop:pdf2-rh-telescopia}
Para cada \(N\geq1\),
\begin{equation}\label{eq:pdf2-rh-telescopia-finita}
 \mathcal B_{{\rm sp},R}^{\rm str}
 =\sum_{n=0}^{N-1}(DA^nE_R)^*(DA^nE_R)
 +(A^NE_R)^*(A^NE_R).
\end{equation}
Equivalentemente, evaluada en \(f,g\in\mathcal D_R\), la identidad finita de
la coligación es
\begin{equation}\label{eq:pdf2-rh-balance-finito-columnas}
\boxed{
\begin{aligned}
 &\langle J_{+,R}f,J_{+,R}g\rangle
  -\langle J_{-,R}f,J_{-,R}g\rangle\\
 &\quad=\sum_{n=0}^{N-1}
  \langle DA^nE_Rf,DA^nE_Rg\rangle
  +\langle A^NE_Rf,A^NE_Rg\rangle.
\end{aligned}}
\end{equation}
El último sumando satisface
\begin{equation}\label{eq:pdf2-rh-terminal-extincion}
 (A^NE_R)^*(A^NE_R)
 =q^{2N}\mathcal B_{{\rm sp},R}^{\rm str}
 \xrightarrow[N\to\infty]{}0.
\end{equation}
\end{proposition}

\begin{proof}
Aplicando \eqref{eq:pdf2-rh-conservacion-local} a \(A^nE_R\),
\[
 (A^nE_R)^*(A^nE_R)
 =(DA^nE_R)^*(DA^nE_R)
 +(A^{n+1}E_R)^*(A^{n+1}E_R).
\]
La suma desde \(n=0\) hasta \(N-1\) cancela sólo los términos intermedios y
conserva \((A^NE_R)^*(A^NE_R)\). Por
\eqref{eq:pdf2-rh-binding-espejo}, \(A^NE_R=q^NE_R\); por tanto el terminal
es \(q^{2N}E_R^*E_R=q^{2N}\mathcal B_{{\rm sp},R}^{\rm str}\). Como
\(0<q=5/9<1\), converge geométricamente a cero.
\end{proof}

Sólo después de demostrar \eqref{eq:pdf2-rh-terminal-extincion} se define
\begin{equation}\label{eq:pdf2-rh-l-r}
 \mathscr L_Rf=(DA^nE_Rf)_{n\geq0}.
\end{equation}
El límite de la telescopía finita da
\begin{equation}\label{eq:pdf2-rh-lstar-l}
 \boxed{
 \mathscr L_R^*\mathscr L_R
 =\mathcal B_{{\rm sp},R}^{\rm str}
 =\mathscr I_{\rm GW}\big|_{\mathcal D_R}\succeq0.}
\end{equation}

\section{Codificador común, proyección y doble publicación explícitas}

Fijemos \(m\geq1\) y dos vectores ortonormales
\(u_m,\eta_m\in\mathbb C^{9^m}\), donde \(u_m\) es el modo uniforme y
\(\eta_m\) el modo de acarreo. El espacio y la proyección del codificador son
\begin{equation}\label{eq:pdf2-rh-hilbert-comun}
 \mathcal X_{m,R}^{\rm sp}
 :=\mathbb C^{9^m}\widehat\otimes
   (H_{-,R}\oplus\mathcal H_R^{\rm sh}),
 \qquad
 P_{m,R}:=|u_m\rangle\langle u_m|\otimes
 I_{H_{-,R}\oplus\mathcal H_R^{\rm sh}},
 \qquad Q_{m,R}:=I-P_{m,R}.
\end{equation}
La acción completa del codificador de los cinco canales es
\begin{equation}\label{eq:pdf2-rh-xi-comun-explicita}
 \boxed{
 \begin{aligned}
  \mathfrak C_{m,R}f
  &:=u_m\otimes(J_{-,R}f,0)
    +\eta_m\otimes(0,E_Rf),\\
  \Xi_{m,R}f&:=\mathfrak C_{m,R}f.
 \end{aligned}}
\end{equation}
No se toma una raíz de \(\mathscr I_{\rm GW}\): las dos entradas son la
salida y el estado de la coligación source--first ya definidos en
\eqref{eq:pdf2-rh-source-state-output}.

\begin{proposition}[Los dos momentos se obtienen de la acción del codificador]
\label{prop:pdf2-rh-vmas-isometria}
Para todo \(f,g\in\mathcal D_R\),
\begin{equation}\label{eq:pdf2-rh-dos-momentos-construidos}
\begin{aligned}
 \langle\mathfrak C_{m,R}f,\mathfrak C_{m,R}g\rangle
  &=\langle J_{+,R}f,J_{+,R}g\rangle,\\
 \langle P_{m,R}\mathfrak C_{m,R}f,
         P_{m,R}\mathfrak C_{m,R}g\rangle
  &=\langle J_{-,R}f,J_{-,R}g\rangle,\\
 \langle Q_{m,R}\Xi_{m,R}f,Q_{m,R}\Xi_{m,R}g\rangle
  &=\mathscr I_{\rm GW}(f,g).
\end{aligned}
\end{equation}
\end{proposition}

\begin{proof}
La ortogonalidad \(u_m\perp\eta_m\) y
\eqref{eq:pdf2-rh-energia-source-first} dan
\[
 \langle\mathfrak C_{m,R}f,\mathfrak C_{m,R}g\rangle
 =\langle J_{-,R}f,J_{-,R}g\rangle
  +\langle E_Rf,E_Rg\rangle
 =\langle J_{+,R}f,J_{+,R}g\rangle.
\]
La proyección \(P_{m,R}\) conserva el sumando uniforme y elimina el de
acarreo; \(Q_{m,R}\) hace lo contrario. Las otras dos identidades siguen de
\eqref{eq:pdf2-rh-binding-source-first}.
\end{proof}

El codificador recién escrito es un representante explícito del lector
estructural inicial, no un lector sustituto. En los subespacios alcanzables se
define
\begin{equation}\label{eq:pdf2-rh-isomorfismo-lector-estructural}
\begin{aligned}
 \mathcal U_{m,R}^{\rm str}
  (P_R^{\rm str}\Xi_R^{\rm str}f)
   &:=u_m\otimes(J_{-,R}f,0),\\
 \mathcal U_{m,R}^{\rm str}
  ((I-P_R^{\rm str})\Xi_R^{\rm str}f)
   &:=\eta_m\otimes(0,E_Rf).
\end{aligned}
\end{equation}
Las dos fuentes son ortogonales. La primera igualdad de Gram de cada fuente
es \eqref{eq:pdf2-rh-momentos-lector-estructural}; para la segunda se usa
\(E_R=\jmath_R^{\rm sh}(I-P_R^{\rm str})\Xi_R^{\rm str}\). Por ello las dos
reglas están bien definidas, son isométricas y se extienden a la suma ortogonal
de sus cierres. En ese cierre alcanzable satisfacen
\begin{equation}\label{eq:pdf2-rh-entrelazamiento-lector-estructural}
 \mathcal U_{m,R}^{\rm str}\Xi_R^{\rm str}=\Xi_{m,R},
 \qquad
 \mathcal U_{m,R}^{\rm str}P_R^{\rm str}
 =P_{m,R}\mathcal U_{m,R}^{\rm str},
 \qquad
 (\Xi_R^{\rm str})^*(I-P_R^{\rm str})\Xi_R^{\rm str}=E_R^*E_R.
\end{equation}
Así, \(B_{{\rm sp},R}^{\rm str}\), el complemento de la estructura de
partida, es exactamente el complemento del codificador construido; no se ha
cambiado de objeto al pasar de \eqref{eq:pdf2-rh-res-sp-completa} a
\eqref{eq:pdf2-rh-xi-comun-explicita}.

La regla
\begin{equation}\label{eq:pdf2-rh-vmas-rango}
 V_{+,m,R}^{0}(J_{+,R}f):=\Xi_{m,R}f
\end{equation}
está bien definida y es isométrica por la primera igualdad de
\eqref{eq:pdf2-rh-dos-momentos-construidos}; se extiende de manera única a
una isometría
\(V_{+,m,R}:H_{+,R}\to\mathcal X_{m,R}^{\rm sp}\). La publicación negativa
es la isometría
\begin{equation}\label{eq:pdf2-rh-vmenos}
 V_{-,m,R}:H_{-,R}\longrightarrow\mathcal X_{m,R}^{\rm sp},
 \qquad V_{-,m,R}y=u_m\otimes(y,0).
\end{equation}
Por construcción,
\begin{equation}\label{eq:pdf2-rh-doble-publicacion-material}
 \boxed{
 V_{+,m,R}J_{+,R}=\Xi_{m,R},
 \qquad
 V_{-,m,R}J_{-,R}=P_{m,R}\Xi_{m,R}.}
\end{equation}
El complemento conserva exactamente el estado antiespecular y su columna de
observación:
\begin{equation}\label{eq:pdf2-rh-complemento-comun}
 \boxed{
 \Xi_{m,R}^*Q_{m,R}\Xi_{m,R}
 =E_R^*E_R
 =\mathscr L_R^*\mathscr L_R
 =\mathscr I_{\rm GW}\big|_{\mathcal D_R}.}
\end{equation}

\begin{proposition}[Conector inducido y coincidencia source--first]
\label{prop:pdf2-rh-conector}
El operador
\begin{equation}\label{eq:pdf2-rh-k-r}
 \mathcal K_R
 :=V_{-,m,R}^*P_{m,R}V_{+,m,R}:
 H_{+,R}\longrightarrow H_{-,R}
\end{equation}
es contractivo, coincide con \(\mathcal K_R^{\rm src}\) en todo
\(H_{+,R}\) y satisface
\begin{equation}\label{eq:pdf2-rh-k-binding}
 \|\mathcal K_R\|\leq1,
 \qquad
 \mathcal K_RJ_{+,R}=J_{-,R}.
\end{equation}
\end{proposition}

\begin{proof}
Las dos cartas son isometrías y \(P_{m,R}\) es una proyección ortogonal,
luego \(\|\mathcal K_R\|\leq1\). Sobre el rango denso de \(J_{+,R}\),
\[
 \mathcal K_RJ_{+,R}f
 =V_{-,m,R}^*P_{m,R}\Xi_{m,R}f
 =V_{-,m,R}^*V_{-,m,R}J_{-,R}f
 =J_{-,R}f.
\]
La misma identidad vale para \(\mathcal K_R^{\rm src}\) por
\eqref{eq:pdf2-rh-output-binding-source}; la continuidad y la densidad
prueban que ambos operadores coinciden.
\end{proof}

La diferencia de Gram queda así publicada en todas sus realizaciones
coincidentes:
\begin{equation}\label{eq:pdf2-rh-tres-publicaciones}
\boxed{
\begin{aligned}
 J_{+,R}^*J_{+,R}-J_{-,R}^*J_{-,R}
 &=\Xi_{m,R}^*Q_{m,R}\Xi_{m,R}\\
 &=E_R^*E_R\\
 &=\mathscr L_R^*\mathscr L_R\\
 &=\mathcal B_{{\rm sp},R}^{\rm str}=\mathscr I_{\rm GW}.
\end{aligned}}
\end{equation}

\section{Naturalidad cofinal y criterio reversible de Weil}

Para \(R\leq R'\), sea
\(i_{R,R'}:\mathcal D_R\hookrightarrow\mathcal D_{R'}\). Las columnas
individuales no forman un sistema isométrico: al aparecer una longitud nueva
\(R<\ell_{p,k}\leq R'\), las dos adquieren la misma energía positiva. Para
\(f,g\in\mathcal D_R\), la separación de soportes da exactamente
\begin{equation}\label{eq:pdf2-rh-naturalidad}
 \left\langle(I-T_{\ell_{p,k}})f,
                    (I-T_{\ell_{p,k}})g\right\rangle
 =2\langle f,g\rangle
 =\left\langle\sqrt2j_mf,\sqrt2j_mg\right\rangle .
\end{equation}
Por tanto el incremento primo de
\(J_{+,R'}^*J_{+,R'}\) coincide con el de
\(J_{-,R'}^*J_{-,R'}\) y se cancela en la diferencia, pero no se declara una
isometría falsa entre las columnas. La compatibilidad cofinal correcta es la
de las formas:
\begin{equation}\label{eq:pdf2-rh-naturalidad-p}
 \boxed{
 \mathscr I_{{\rm GW},R}
 =i_{R,R'}^*\mathscr I_{{\rm GW},R'}i_{R,R'},
 \qquad
 \mathcal B_{{\rm sp},R}^{\rm str}
 =i_{R,R'}^*\mathcal B_{{\rm sp},R'}^{\rm str}i_{R,R'}.}
\end{equation}
Como \(E_R^*E_R=\mathscr L_R^*\mathscr L_R=\mathscr I_{{\rm GW},R}\), las
realizaciones mínimas de Kolmogórov sí admiten isometrías únicas sobre sus
rangos alcanzables, definidas por
\begin{equation}\label{eq:pdf2-rh-naturalidad-bw}
 \iota_{R,R'}^E(E_Rf):=E_{R'}f,
 \qquad
 \iota_{R,R'}^L(\mathscr L_Rf):=\mathscr L_{R'}f
 \quad(f\in\mathcal D_R).
\end{equation}
Estas dos reglas están bien definidas precisamente por
\eqref{eq:pdf2-rh-naturalidad-p}; no se extienden por decreto a
\(J_{\pm,R}\), \(\Xi_{m,R}\) o \(\mathcal K_R\). El refinamiento nonádico
\(m\mapsto m+1\) permanece, a ventana fija, en las isometrías estructurales
que transportan \(j_m\), \(\widehat\delta_{a,m}\), \(u_m\) y \(\eta_m\).

Para no comprimir el último paso, fijamos también sus objetos y
normalizaciones. Sea
\[
 \xi(s):=\frac12s(s-1)\pi^{-s/2}\Gamma(s/2)\zeta(s)
\]
y sea \(\mathscr Z\) el multiconjunto de sus ceros, cada cero repetido según
su multiplicidad. La ecuación funcional y la realidad de los coeficientes
preservan \(\mathscr Z\) por
\(\rho\mapsto1-\rho\), \(\rho\mapsto\overline\rho\) y
\(\rho\mapsto1-\overline\rho\). Con nuestra convención de Fourier se define
\begin{equation}\label{eq:pdf2-rh-transformada-weil}
 \Phi_f(s):=\widehat f\!\left(i(s-\tfrac12)\right),
 \qquad
 \widetilde f(x):=\overline{f(-x)}.
\end{equation}
La fórmula explícita completa calculada en
\eqref{eq:pdf2-rh-gram-difference-gw}, incluidos los términos gamma y polar
que retiran polos y ceros triviales, tiene la forma espectral
\begin{equation}\label{eq:pdf2-rh-formula-espectral-weil}
 \mathscr I_{\rm GW}(f,g)
 =\sum_{\rho\in\mathscr Z}
   \Phi_f(\rho)\,
   \overline{\Phi_g(1-\overline\rho)}.
\end{equation}
La suma cuenta multiplicidades. Converge absolutamente: para
\(f,g\in C_c^\infty(\mathbb R)\), Paley--Wiener da decaimiento más rápido
que cualquier potencia en bandas verticales, mientras el número de ceros con
\(|\operatorname{Im}\rho|\leq T\) es \(O(T\log T)\).

\begin{proposition}[Criterio reversible de Weil con cuantificadores completos]
\label{prop:pdf2-rh-criterio-weil-desplegado}
Bajo la normalización
\eqref{eq:pdf2-rh-transformada-weil} y contando todos los ceros no triviales
con su multiplicidad,
\begin{equation}\label{eq:pdf2-rh-criterio-weil}
 \left[
  \mathscr I_{\rm GW}(f,f)\geq0
  \text{ para todo }f\in C_c^\infty(\mathbb R;\mathbb C)
 \right]
 \Longleftrightarrow
 \left[
  \operatorname{Re}\rho=\frac12
  \text{ para todo }\rho\in\mathscr Z
 \right].
\end{equation}
\end{proposition}

La implicación recíproca no se obtiene interpolando un número creciente de
ceros: esa operación no daría control uniforme sobre la cola. Se usa el
criterio clásico completo de Weil en la formulación de Bombieri, cuyo espacio
de prueba y cuya normalización se trasladan aquí sin abreviarlos.\footnote{%
E.~Bombieri, \emph{Remarks on Weil's quadratic functional in the theory of
prime numbers, I}, Rendiconti Lincei, Matematica e Applicazioni, ser.~9,
vol.~11 (2000), pp.~183--233. El criterio allí fijado afirma que el funcional
cuadrático de la fórmula explícita es semidefinido positivo sobre toda
\(C_c^\infty(0,\infty)\) si y sólo si vale la hipótesis de Riemann.}

\begin{proof}
Escribimos materialmente el cambio de coordenadas que identifica ambos
criterios. Para \(f\in C_c^\infty(\mathbb R)\), sea
\begin{equation}\label{eq:pdf2-rh-cambio-mellin}
 (\mathcal M f)(x):=x^{-1/2}f(\log x),
 \qquad x>0.
\end{equation}
El mapa
\[
 \mathcal M:C_c^\infty(\mathbb R)
 \longrightarrow C_c^\infty(0,\infty)
\]
es biyectivo, con inversa
\begin{equation}\label{eq:pdf2-rh-cambio-mellin-inverso}
 (\mathcal M^{-1}g)(u)=e^{u/2}g(e^u).
\end{equation}
Si
\(\widehat g_{\rm Mel}(s):=\int_0^\infty g(x)x^{s-1}\,dx\), el cambio
\(x=e^u\) da, sin factor residual,
\begin{equation}\label{eq:pdf2-rh-mellin-fourier}
 \begin{aligned}
 \widehat{\mathcal M f}_{\rm Mel}(s)
 &=\int_{\mathbb R}f(u)e^{(s-1/2)u}\,du\\
 &=\widehat f\!\left(i(s-\tfrac12)\right)
 =\Phi_f(s).
 \end{aligned}
\end{equation}
Además,
\begin{equation}\label{eq:pdf2-rh-mellin-involucion}
 \widehat{\overline{\mathcal M f}}_{\rm Mel}(1-\rho)
 =\overline{
   \widehat{\mathcal M f}_{\rm Mel}(1-\overline\rho)}
 =\overline{\Phi_f(1-\overline\rho)}.
\end{equation}
Por tanto el funcional del criterio de Weil--Bombieri es exactamente
\begin{equation}\label{eq:pdf2-rh-bombieri-identificado}
 \sum_{\rho\in\mathscr Z}
 \widehat{\mathcal M f}_{\rm Mel}(\rho)\,
 \widehat{\overline{\mathcal M f}}_{\rm Mel}(1-\rho)
 =
 \sum_{\rho\in\mathscr Z}
 \Phi_f(\rho)\overline{\Phi_f(1-\overline\rho)}
 =\mathscr I_{\rm GW}(f,f).
\end{equation}
No se cambia la clase de pruebas, no se pierde ninguna multiplicidad y no se
trunca el conjunto de ceros, porque \(\mathcal M\) es una biyección. El
criterio clásico citado aplicado a
\eqref{eq:pdf2-rh-bombieri-identificado} da precisamente la equivalencia
\eqref{eq:pdf2-rh-criterio-weil}. En la dirección directa, si
\(\operatorname{Re}\rho=1/2\), la misma identidad se reduce explícitamente a
\[
 \mathscr I_{\rm GW}(f,f)
 =\sum_{\rho\in\mathscr Z}|\Phi_f(\rho)|^2\geq0.
\]
\end{proof}

La igualdad \(\bigcup_R\mathcal D_R=C_c^\infty(\mathbb R)\) es exacta; cada
función empleada en el argumento pertenece a alguna ventana y la naturalidad
cofinal transporta allí la positividad ya demostrada.

\begin{theorem}[Resolución Riemann--Weil desde la imagen espectral--polar]
\label{thm:pdf2-rh-resolucion}
Fijada la estructura completa
\(\mathfrak G_{\rm sp}\), la forma
completa de Guinand--Weil es positiva en la clase cofinal y separadora. En
consecuencia, todo cero no trivial \(\rho\) de la función zeta de Riemann
satisface
\[
 \boxed{\operatorname{Re}\rho=\frac12.}
\]
\end{theorem}

\begin{proof}
La expansión primo--gamma--polar identifica la diferencia de las dos columnas
con \(\mathscr I_{\rm GW}\) usando un único regulador y lo retira por
convergencia dominada. La recurrencia antiespecular
\eqref{eq:pdf2-rh-recurrencia-residual} anula la diferencia de salida antes de
formar una desigualdad; por tanto la isometría \(\Sigma_R\) da
\eqref{eq:pdf2-rh-energia-source-first}. Comparar esa identidad con la
expansión calculada produce
\(E_R^*E_R=\mathscr I_{\rm GW}\). La identidad local
\(A^*A+D^*D=I\) se itera sin borrar el terminal y da
\eqref{eq:pdf2-rh-telescopia-finita}; el terminal es exactamente
\(q^{2N}E_R^*E_R\) y se extingue sólo en el límite. De ahí
\(\mathscr L_R^*\mathscr L_R=\mathscr I_{\rm GW}\succeq0\).
La acción
\eqref{eq:pdf2-rh-xi-comun-explicita} construye entonces el codificador, sus
dos momentos, la proyección y las dos publicaciones, y
\eqref{eq:pdf2-rh-tres-publicaciones} identifica todas las realizaciones del
mismo cuadrado. La naturalidad transporta esa identidad a toda la unión
cofinal. Finalmente,
\eqref{eq:pdf2-rh-criterio-weil} concluye la afirmación.
\end{proof}

\section{Control Redheffer local posterior y no gobernante}

El siguiente despliegue verifica que la pasividad local puede obtenerse desde
el acarreo sin usar la forma de Weil. Es un control de no-circularidad y no una
premisa del \cref{thm:pdf2-rh-resolucion}. Esta transferencia truncada no es
la coligación completa \(\Sigma_R\) de
\eqref{eq:pdf2-rh-coligacion-source-first}.

Para \(N=9^m\), sea \(U_{a,N}\) el odómetro unitario
\[
 U_{a,N}(e_j\otimes f)=e_{j+1}\otimes f\ (j<N-1),
 \qquad U_{a,N}(e_{N-1}\otimes f)=e_0\otimes T_af.
\]
Con \(P_N=|u_N\rangle\langle u_N|\otimes I\), \(Q_N=I-P_N\),
\(A_{a,N}=P_NU_{a,N}P_N\) y \(C_{a,N}=Q_NU_{a,N}P_N\),
\begin{equation}\label{eq:pdf2-rh-source-energia}
 A_{a,N}^*A_{a,N}+C_{a,N}^*C_{a,N}=P_N,
 \qquad
 C_{a,N}^*C_{a,N}
 =\frac{N-1}{N^2}(2I-T_a-T_a^*).
\end{equation}
Para \(0<s<1\), la transferencia
\begin{equation}\label{eq:pdf2-rh-source-redheffer}
 \mathscr R_{a,N}(s)=(sI-A_{a,N})(I-sA_{a,N})^{-1}
\end{equation}
satisface
\begin{equation}\label{eq:pdf2-rh-source-defecto}
\begin{aligned}
 I-\mathscr R_{a,N}(s)^*\mathscr R_{a,N}(s)
 ={}&(1-s^2)(I-sA_{a,N}^*)^{-1}\\
 &\quad\cdot C_{a,N}^*C_{a,N}(I-sA_{a,N})^{-1}\succeq0.
\end{aligned}
\end{equation}
Con
\[
 s_{p,N}=\frac{Np^{-1/2}}{1+(N-1)p^{-1/2}},
 \qquad s_\Gamma(a)=e^{-a/2},
\]
la suma sobre relojes primos y la integral directa gamma conservan norma a lo
sumo uno. De forma explícita,
\begin{equation}\label{eq:pdf2-rh-source-relojes}
 \mathscr R_{N,R}
 =\bigoplus_{\log p\leq R}
   \mathscr R_{\log p,N}(s_{p,N})
 \ \oplus\!
 \int_0^\infty{}^\oplus
   \mathscr R_{a,N}(s_\Gamma(a))\,da,
 \qquad \|\mathscr R_{N,R}\|\leq1.
\end{equation}
La carta finita de roles
\[
 \mathcal H_+=\operatorname{span}\{f_3,f_6,f_9\},\quad
 \mathcal H_-=\operatorname{span}\{f_4,f_8\},\quad
 C_\angle=|f_4\rangle\langle f_3|+|f_8\rangle\langle f_9|
\]
satisface \(I-C_\angle^*C_\angle=P_6\). Con las cartas energéticas
two-way se fija
\begin{equation}\label{eq:pdf2-rh-source-cartas}
\begin{aligned}
 \mathcal K_q&=K_2\otimes(\mathbb C^9)^{\otimes q},\\
 \mathsf B_{\pm,q}:\mathcal M_{\pm,q,R}
 &\longrightarrow
 \mathcal K_q\widehat\otimes\mathcal H_\pm
 \widehat\otimes\mathcal H_{\mathfrak C_R},\\
 \mathsf B_{\pm,q}^*\mathsf B_{\pm,q}&=\mathsf G_{\pm,q},
 \qquad
 \mathsf B_{\pm,q}^{\sharp}
 =\mathsf G_{\pm,q}^{-1}\mathsf B_{\pm,q}^*.
\end{aligned}
\end{equation}
La transferencia
\begin{equation}\label{eq:pdf2-rh-source-transferencia}
 \mathscr C_{q,R}^{\rm src}
 =\mathsf B_{-,q}^{\sharp}
  (I\otimes C_\angle\otimes\mathscr R_{N,R})\mathsf B_{+,q}
\end{equation}
tiene defecto
\begin{align}
 \Delta_{q,R}^{\rm src}
 ={}&I\otimes\bigl[P_6\otimes I
 +(P_3+P_9)\otimes
 (I-\mathscr R_{N,R}^*\mathscr R_{N,R})\bigr]\succeq0,
 \label{eq:pdf2-rh-source-delta}\\
 I-(\mathscr C_{q,R}^{\rm src})^\sharp\mathscr C_{q,R}^{\rm src}
 ={}&\mathsf B_{+,q}^{\sharp}\Delta_{q,R}^{\rm src}
 \mathsf B_{+,q}\succeq0.
 \label{eq:pdf2-rh-source-contractividad}
\end{align}
Este control conserva ruta, reloj, rol, acarreo y los canales antes de formar
una diferencia. No se identifica con \(\mathcal K_R\) y no sustituye la
publicación común explícita de
\eqref{eq:pdf2-rh-doble-publicacion-material}; por ello su alcance local no
reabre la conclusión cofinal.

\section{Falsadores de una sombra ablacionada}

\begin{falsifier}[Sustituciones que cambian el objeto de partida]
El argumento anterior no se transfiere a una sombra que elimine una de las
trece coordenadas de \eqref{eq:pdf2-rh-res-sp-completa}, construya \(J_+\) y
\(J_-\) desde antecedentes distintos, reemplace una proyección ortogonal por
un idempotente no autoadjunto, omita un canal, use reguladores incompatibles,
escoja \(P_R\) después del signo, borre el terminal antes del límite o reduzca
la clase a una familia no cofinal.
\end{falsifier}

\begin{proof}
Cada operación rompe, respectivamente, el dominio de
\(\mathfrak G_{\rm sp}\), la buena definición de
\eqref{eq:pdf2-rh-vmas-rango}, el cuadrado ortogonal, la identidad funcional,
la dirección causal, la telescopía finita o el criterio reversible. El
falsador afecta a la sombra sustituida; no demueve la imagen estructural
completa empleada en el \cref{thm:pdf2-rh-resolucion}.
\end{proof}

\paragraph{Procedencia y fuerza probatoria.}
La restricción espectral--polar, la coligación source--first, su complemento y
la hoja espejo son \texttt{ARQUITECTURA\_AUTORAL\_PREEXISTENTE}. La
anulación coinductiva de \(\Omega\), la identidad finita con terminal y el
codificador de acción
\eqref{eq:pdf2-rh-xi-comun-explicita} son
\texttt{FORMALIZACION\_NUEVA} de esa estructura completa. La identificación
Guinand--Weil y el cierre cofinal son \texttt{RESULTADO\_RECUPERADO}. Las
ecuaciones
\eqref{eq:pdf2-rh-coligacion-source-first}--
\eqref{eq:pdf2-rh-tres-publicaciones} exhiben el objeto, sus dominios, su
acción, el terminal y el límite sin introducir una raíz abstracta de la forma;
su fuerza es \emph{demostrado con estructura de partida explícita}.
